\section{The global DARC compilation theorem}
\label{sec:global}

We now assemble the preceding components into one conditional theorem.  The
statement is deliberately explicit about the boundary between generic
compilation and theorem-specific mathematics.

\subsection{Well-formed plans}

A normalized fixed plan \(P\) is \emph{well formed} when the following data
and proofs are supplied.

\begin{enumerate}
  \item A finite resource shape \(L_P\), a query seed \(q_P\), and a monotone
        transfer \(F_P:L_P\to L_P\).
  \item A finite proof graph whose condensation after explicit loop extraction
        is acyclic.
  \item For each demand state \(d\) and each node active at \(d\), an ordinary
        local contract from active dependencies and selected boundary formulas
        to the node conclusion.
  \item For each loop node, concrete trace symbols, defining equations, a
        quantifier-free invariant, and base/step derivations.
  \item For each bounded-formula node, reflected atoms and the generated
        fixed-formula reflection derivation.
  \item For each relation node, finite bounds and reflected edge/terminal bits.
  \item For each scoped certificate, a local producer and consumer, or an exact
        producer formula placed in the residual boundary.
  \item For each named certificate, a finite bound, validity bit, reflection
        theorem, and a proof that some valid code lies below the bound.
  \item For each quantitative operation, a sound transfer theorem connecting
        requested output resources to input precisions, guards, and residuals.
  \item For each backend export, checked vocabulary comparisons and a Bridge
        theorem preserving the required open or closed sequent.
\end{enumerate}

The plan may contain open residual producers.  Well-formedness requires that
they be named by exact formulas; it does not require derivations of them.

\subsection{Compiler output}

Compute \(d^*=\cl_P(q_P)\) by the finite iteration of
\cref{sec:demand}.  Define:
\begin{itemize}
  \item \(V^*(P)\), the active nodes at \(d^*\);
  \item \(B^*(P)\), the residual and guard formulas selected at \(d^*\) that
        are not already discharged;
  \item \(\Phi_P\), the public conclusion specialized to the requested
        outputs and precisions;
  \item \(S^*(P)\), a topological schedule of the retained condensation graph,
        with loop, bounded-formula, relation, and certificate blocks expanded
        through their dedicated compilers.
\end{itemize}
The compiler's primary output type is
\[
  B^*(P)\vdash_{\PRAAx}\Phi_P.
  \tag{13.1}\label{eq:compiled-endpoint}
\]
If the query has free parameters, both context and conclusion remain in that
finite scope.

\begin{theorem}[DARC soundness]
\label{thm:darc-soundness}
Every well-formed fixed plan \(P\) has a finite compilation that constructs an
ordinary derivation of \eqref{eq:compiled-endpoint}.
\end{theorem}

\begin{proof}
The proof has five stages.

\paragraph{1. Demand closure.}
By \cref{thm:demand-stabilization}, computation of \(d^*\) terminates.  By
\cref{thm:demand-least}, it contains every resource forced by the query and
local transfer summaries.  By transfer soundness, all resources used by a
retained local proof are present.

\paragraph{2. Specialized block proofs.}
For each retained bounded-formula node, invoke the generated reflection theorem
\cref{thm:bounded-reflection}.  For each retained loop node, invoke normalized
trace reflection \cref{thm:trace-reflection}.  For each relation node, invoke
finite viability and, only if needed, canonical finite path construction
\cref{thm:finite-viability,thm:least-path}.  These operations produce ordinary
local sequents.

\paragraph{3. Certificate handling.}
For every sealed certificate, compose its producer and consumer by existential
elimination \cref{thm:local-cert-erasure}.  For an observed certificate, insert
the supplied compatibility theorem and seal the witness.  For a named
certificate, invoke bounded least search \cref{thm:named-export}.  If a
producer is absent, keep its exact existential or quantitative formula in
\(B^*(P)\).

\paragraph{4. Acyclic proof assembly.}
After block expansion, the retained condensation graph is acyclic.  Compile it
in topological order by \cref{thm:acyclic-compile}, using simultaneous cut to
supply dependencies.  Every selected boundary formula is available either by
the hypothesis rule or by a producer derivation already inserted into the
schedule.

\paragraph{5. Output assembly.}
Combine retained public outputs by the ordinary logical introduction rules and
specialize them at the requested terms.  The result is a finite inhabitant of
\(B^*(P)\vdash_{\PRAAx}\Phi_P\).
\end{proof}

\subsection{Residual exactness}

Soundness alone would allow the compiler to retain every conceivable
hypothesis.  Demand leastness gives a stronger audit property.

\begin{theorem}[Residual exactness relative to summaries]
\label{thm:residual-exactness}
Let \(e\) be any demand state above the query that satisfies all declared local
transfer summaries.  Then \(d^*\le e\).  Consequently every residual, guard,
precision, or certificate coordinate selected by DARC is forced by the
published summaries.
\end{theorem}

\begin{proof}
This is \cref{thm:demand-least}.  The interpretation of the coordinates gives
the consequence.
\end{proof}

The theorem does not say that every selected residual is logically necessary
for \(\Phi_P\).  It says that it is necessary for this modular proof plan with
these local contracts.  A better local theorem can legitimately improve the
endpoint.

\subsection{Producer substitution and monotone improvement}

Let \(B^*(P)=[B_1,\ldots,B_r]\).  If later work supplies derivations of some
\(B_i\), \cref{cor:incremental-discharge} removes them.  If a supplied producer
also has a better transfer summary, rerunning demand analysis may remove
additional nodes and boundaries.

\begin{corollary}[Full discharge]
If \(\vdash B_i\) for every residual in \(B^*(P)\), then
\[
  \vdash_{\PRAAx}\Phi_P.
\]
\end{corollary}

This is the preferred route from research packet to checked theorem: first
compile an exact conditional endpoint, then discharge producers independently.

\subsection{Query monotonicity}

If query seed \(q\le q'\), monotonicity of the update implies
\[
  \cl_P(q)\le\cl_P(q').
\]
Thus asking for more outputs cannot reduce the declared resource demand.  A
value-only query may hide a branch witness; a value-plus-witness query can only
add the branch, separation guard, and supporting proof nodes.

\begin{proof}
Induct simultaneously over the inflationary iterations from \(q\) and \(q'\).
The base inequality is given.  Monotonicity of \(F_P\) and of join preserves it
at every stage.  Pass to stable values.
\end{proof}

\subsection{Plan composition}

Let \(P\) compile a conclusion used as an input producer of \(Q\).  Form the
composite by identifying that interface formula, extending each resource shape
with zeros for the other component, and joining transfer maps.  The global
compiler may inline both schedules or retain the theorem boundary and use
simultaneous cut.  The latter is preferable when \(P\) is a stable library
component.

\begin{proposition}[Compositional soundness]
If \(P\) and \(Q\) are well formed and their shared formulas and resource
interfaces agree, then the composite plan is well formed and its compiled
endpoint is obtained by ordinary cut composition of the two endpoints,
possibly after a smaller joint demand closure prunes unused outputs of \(P\).
\end{proposition}

\begin{proof}
Monotonicity is preserved by extending and joining the transfers.  The
combined graph is finite; any cross-component cycle is either a resource cycle,
handled by demand closure, or must be normalized into an explicit loop block.
Local contracts remain valid under weakening and substitution.  Compile the
joint retained graph by \cref{thm:darc-soundness}; at the theorem boundary the
composition is \cref{lem:simcut}.
\end{proof}

\subsection{Backend preservation}

After obtaining \eqref{eq:compiled-endpoint}, an application may request a
Foundation sequent.  The checked open Bridge translates the context and
conclusion under an explicit assignment.  Symbol-comparison theorems rewrite
the generic arithmetic vocabulary to the client's existing codes.  DARC treats
this as a final node so that backend requirements do not contaminate the
mathematical core.

No theorem is identified merely because two Lean functions compute the same
standard natural-number function.  The bridge requires an object-level
comparison or an exact syntax round trip.  This preserves the theory and
encoding stated by each endpoint.
